\InterpolationArticle{7}
Der in der Praxis häufigste Fall liegt vor, wenn \(a,b,c,d\mathinner{.\,.\,.}\) eine arithmetische Folge bilden. Setzt man den Abstand \(=1\), sodass \(b=a+1,\ c=a+2\)
usw. gilt, so wird die Formel von Artikel 4 zu
\[
\begin{aligned}
T={}&A\\
&+(t-1)(B-A)\\
&+(t-1)(t-2)\left(\frac{C-2B+A}{2}\right)\\
&+(t-1)(t-2)(t-3)\left(\frac{D-3C+3B-A}{6}\right)\\
&+\text{usw.}
\end{aligned}
\]
oder
\[\begin{aligned}
T&=A+A'(t-1)+A''(t-1)(t-2)\\&+A'''(t-1)(t-2)(t-3)+\mathinner{.\,.}
\end{aligned}\]
wobei \(A',A'',A'''\) usw. nach dem folgenden Algorithmus berechnet werden
\[
\begin{array}{lll}
A'=B-A,&B'=C-B,&C'=D-C\quad\text{usw.}\\
2A''=B'-A',&2B''=C'-B',&2C''=D'-C'\quad\text{usw.}\\
3A'''=B''-A'',&3B'''=C''-B'',&3C'''=D''-C''\quad\text{usw.}
\end{array}
\]
usw. Diese Formel stimmt mit der üblichen Interpolationsformel durch Differenzen vollständig überein.

\Needspace{42\baselineskip}
\InterpolationArticle{8}
Sind für genügend viele Werte von \(x\), die in der arithmetischen Folge
\(a,a+1,a+2\mathinner{.\,.\,.}\) fortschreiten, die entsprechenden Werte der Funktion \(X\) bekannt, sodass eine Folge von \(m\)
aufeinanderfolgenden Werten von \(x\) nach Belieben zur Berechnung von \(T\) gewählt werden kann:
so entsteht die Frage, \emph{welche} Werte für diesen Zweck am besten vorzuziehen sind, wenn wir weder mehr als \(m\) verwenden noch über die \(m-1^{\mathrm{te}}\) Differenz hinausgehen wollen?
\SourceLeaf{286}
Offenbar ist dies so zu entscheiden, dass der aus der Differenz \(m^{\mathrm{ten}}\) Ordnung entstehende Fehler
möglichst klein wird; dieser Fehler ist
\[
=(t-a)(t-a-1)(t-a-2)\mathinner{.\,.\,.\,.}(t-a-m+1)\frac{\Delta}{1.2.3\mathinner{.\,.\,.} m}
\]
wenn die zu \(x=a,a+1,a+2,\mathinner{.\,.\,.} a+m-1\) gehörigen Glieder verwendet werden, oder
\[
=(t-a-1)(t-a-2)(t-a-3)\mathinner{.\,.\,.\,.}(t-a-m)\frac{\Delta}{1.2.3\mathinner{.\,.\,.} m}
\]
wenn die zu \(x=a+1,a+2,a+3,\mathinner{.\,.\,.} a+m\) gehörigen Glieder verwendet werden, oder
\[
=(t-a+1)(t-a)(t-a-1)\mathinner{.\,.\,.}(t-a-m+2)\frac{\Delta}{1.2.3\mathinner{.\,.\,.\,.}m}
\]
wenn die zu \(x=a-1,a,a+1,\mathinner{.\,.\,.} a+m-2\) gehörigen Glieder verwendet werden, wobei
\(\Delta\) die Differenz \(m^{\mathrm{ten}}\) Ordnung bezeichnet, die hier, da die Differenzen höherer Ordnungen
unberücksichtigt bleiben, als konstant angesehen werden kann. Sind also die \(m\) Glieder,
die zu \(x=a,a+1,a+2,\mathinner{.\,.\,.} a+m-1\) gehören, die geeignetsten, so muss von jenen
drei Ausdrücken der erste der kleinste sein und somit ohne Rücksicht auf das Vorzeichen \(t-a<\)
oder wenigstens nicht \(>t-a-m\) gelten, ebenso \(t-a+1>\) oder wenigstens nicht \(<t-a-m+1\).
Aus der ersten Bedingung folgt leicht, dass \(t-a-m\) eine negative Größe sein muss, die zwischen \(-\frac12m\) und \(-\infty\) liegt; daher liegt \(t-a\) zwischen \(+\frac12m\) und \(-\infty\). Aus
der zweiten Bedingung dagegen folgt, dass \(t-a+1\) eine positive Größe ist und zwischen \(\frac12m\) und
\(+\infty\) liegt; folglich liegt \(t-a\) zwischen \(\frac12m-1\) und \(\frac12m\). Hieraus folgt: Ist \(m\) eine gerade Zahl, so sind die zur Interpolation zu verwendenden Werte so zu wählen, dass
die erste Hälfte auf der einen, die zweite auf der anderen Seite des gesuchten Gliedes liegt. Ist dagegen
\(m\) ungerade, so sind diejenigen Glieder zu wählen, deren mittleres dem gesuchten möglichst nahe liegt.
Liegt in diesem Fall \(t\) genau in der Mitte zwischen zwei aufeinanderfolgenden Werten von \(x\), so macht es hinsichtlich des Fehlers aus der vernachlässigten Differenz \(m^{\mathrm{ten}}\) Ordnung im Allgemeinen keinen Unterschied, ob das dem gesuchten vorangehende oder das ihm folgende Glied als mittleres
der zu verwendenden Glieder gewählt wird.

\InterpolationArticle{9}
\emph{Beispiele.} I. Es sei \(\log.\sin 24^{0}30'\) zu bestimmen, wenn die Logarithmen der Sinus für jeden einzelnen Grad vorliegen, wobei die vierten und höheren Differenzen vernachlässigt werden. In
diesem Fall verwenden wir also nach den Regeln des vorhergehenden Artikels die Logarithmen der Sinus
der Bögen \(23^{0},24^{0},25^{0},26^{0}\). Bezeichnen wir diese mit \(A,B,C,D\), so ergibt sich nach der Formel von Artikel 3 der gesuchte Logarithmus
\SourceLeaf{287}
\[
=-\frac{1}{16}A+\frac{9}{16}B+\frac{9}{16}C-\frac{1}{16}D
\]
\Needspace{11\baselineskip}
dessen Zahlenwert sich aus
\[
\begin{aligned}
A&=9{,}5918780\\
B&=9{,}6093133\\
C&=9{,}6259483\\
D&=9{,}6418420
\end{aligned}
\]
als \(9{,}6177271.5\) ergibt; die Tafeln geben \(9{,}6177270\).

Hätten wir hier anstelle von \(\log.\sin 23^{0}\) den Logarithmus von \(\sin 27^{0}=E=9{,}6570468\) verwendet, so hätte die Formel
\[
+\frac{5}{16}B+\frac{15}{16}C-\frac{5}{16}D+\frac{1}{16}E
\]
den Wert \(9{,}6177267\,.\,375\) ergeben.

Die Formel des vorhergehenden Artikels gibt für diesen Fall den aus der vernachlässigten vierten Differenz
\(\Delta\) entstehenden Fehler als \(=\frac{3}{128}\Delta\), wenn \(A,B,C,D\), dagegen als \(=-\frac{5}{128}\Delta\), wenn \(B,C,D\)\raisebox{-0.3ex}{:}\,\(E\) verwendet werden; dieser Fehler ist zum berechneten Wert zu addieren. In unserem Beispiel gilt
\(\Delta=-0{,}0000066\); die daraus berechnete Korrektur bringt beide Rechnungen mit den Tafeln in Übereinstimmung.

II. Es sei die Zeit des Neumonds am 16. Juni 1806 aus den Längen der Sonne und des Mondes für
die einzelnen Tage zu berechnen, wobei die Differenzen bis einschließlich zur vierten
zu berücksichtigen sind. In den Pariser Ephemeriden stehen für den wahren Mittag des Ortes die folgenden Werte

\begin{center}
\begin{tabular}{r@{\qquad}rrr@{\qquad}rrr}
&\multicolumn{3}{c}{Sonnenlänge}&\multicolumn{3}{c}{Mondlänge}\\
1806 Juni 14&\(82^{0}\)&\(39'\)&\(12''\)&\(53^{0}\)&\(19'\)&\(15''\)\\
15&83&36&30&67&29&43\\
16&84&33&48&81&59&52\\
17&85&31&7&96&44&9\\
18&86&28&24&111&35&30
\end{tabular}
\end{center}

Hier sind die gegebenen Werte von \(x\) die Differenzen zwischen der Sonnen- und der Mondlänge, die entsprechenden Werte von \(X\) die Zeiten, die vom Mittag des 16. Juni an gezählt werden; gesucht ist der Wert von \(X\), der dem Wert \(x=0\) entspricht. Somit gilt
\SourceLeaf{288}
\[
\begin{array}{r@{\;}c@{\;}r@{\;}r@{\;}r@{\quad}c@{\;}r@{\qquad}r@{\;}c@{\;}r}
a&=&-29^{0}&19'&57''&=&-105596''&A&=&-2\\
b&=&-16&6&47&=&-58007&B&=&-1\\
c&=&-2&33&56&=&-9236&C&=&0\\
d&=&+11&13&2&=&+40382&D&=&+1\\
e&=&+25&7&6&=&+90426&E&=&+2\\
t&=&0&&&&&&&
\end{array}
\]
\Needspace{23\baselineskip}
Hieraus wird berechnet
\[
\begin{aligned}
\frac{(t-b)(t-c)(t-d)(t-e)}{(a-b)(a-c)(a-d)(a-e)}&=+0{,}0149084\\[5pt]
\frac{(t-a)(t-c)(t-d)(t-e)}{(b-a)(b-c)(b-d)(b-e)}&=-0{,}1050660\\[5pt]
\frac{(t-a)(t-b)(t-d)(t-e)}{(c-a)(c-b)(c-d)(c-e)}&=+0{,}9624567\\[5pt]
\frac{(t-a)(t-b)(t-c)(t-e)}{(d-a)(d-b)(d-c)(d-e)}&=+0{,}1434435\\[5pt]
\frac{(t-a)(t-b)(t-c)(t-d)}{(e-a)(e-b)(e-c)(e-d)}&=-0{,}0157429
\end{aligned}
\]
Werden diese Koeffizienten\HistoricalNote{23} der Reihe nach mit \(-2,-1,0,1,2\) multipliziert, so ist die Summe der Produkte
\(=+0{,}1872069\); der Neumond ist daher am 16. Juni um \(4^{\mathrm h}29'34''7\) wahrer Pariser
Zeit oder \(4^{\mathrm h}29'41''4\) mittlerer Zeit; die Connaissance des tems gibt \(4^{\mathrm h}27'42''\).

